\chapter{L'ambiente di lavoro}

Prima di scrivere il primo programma bisogna preparare il proprio
\textit{laboratorio}: qualche programma da installare e un paio di abitudini
da prendere. Questo capitolo accompagna passo passo dall'installazione
alla prima esecuzione, su \textbf{Windows}, \textbf{macOS} e \textbf{Linux}.
Ogni passo è spiegato: non serve sapere nulla di computer oltre ad
accenderli e a scaricare un file da Internet.

\warning{Leggi solo la parte del tuo sistema operativo}{
  Le istruzioni di installazione sono divise per sistema operativo. Segui
  quelle del tuo computer e salta le altre: dalla sezione \textit{Il
  terminale} in poi le istruzioni tornano a essere le stesse per tutti.}

\section{Che cosa ci serve}

Un computer non capisce il C. Capisce soltanto il proprio
\term{linguaggio macchina}: lunghe sequenze di numeri che il processore sa
eseguire, una dopo l'altra. Noi invece scriveremo i programmi in C, un
linguaggio fatto di parole inglesi e simboli, comprensibile a una persona.
Fra noi e il processore servono quindi tre strumenti, che compaiono nella
figura~\ref{fig:catena}.

\begin{figure}[H]
  \centering
  \begin{tikzpicture}[
      card/.style={rounded corners=4pt, draw=#1!65!black, fill=#1!6!white,
        line width=0.8pt, minimum width=4.6cm, minimum height=3.3cm},
      titolo/.style={font=\Lato\small\bfseries, text=#1!55!black},
      schermo/.style={fill=vscbg, rounded corners=2pt, inner sep=5pt,
        font=\FiraCodeNL\scriptsize, text=vscfg, align=left},
      nota/.style={font=\Lato\scriptsize, text=neutral!25!black, align=center}]
    \node[card=primary] (a) at (0,0) {};
    \node[card=accent] (b) at (6.6,0) {};
    \node[card=primary] (c) at (13.2,0) {};
    \node[titolo=primary, anchor=north] at ([yshift=-5pt]a.north) {1. Scrivere};
    \node[titolo=accent, anchor=north] at ([yshift=-5pt]b.north) {2. Tradurre};
    \node[titolo=primary, anchor=north] at ([yshift=-5pt]c.north) {3. Eseguire};
    \node[schermo] at ([yshift=-1pt]a.center) {%
      {\color{vscblue}int} {\color{vscfunc}main}({\color{vscblue}void})\\
      \{\\
      \ \ {\color{vscfunc}printf}({\color{vscstring}"Ciao!\textbackslash n"});\\
      \}};
    \node[schermo, text=vscnum] at ([yshift=-1pt]b.center) {%
      01010101 01001000\\10001001 11100101\\01001000 10001101\\11000011 \dots};
    \node[schermo] at ([yshift=-1pt]c.center) {%
      {\color{white}\bfseries\$ ./ciao}\\Ciao!\\[2pt]\phantom{x}};
    \node[nota, anchor=south] at ([yshift=5pt]a.south) {\FiraCodeNL ciao.c\Lato: un file di testo};
    \node[nota, anchor=south] at ([yshift=5pt]b.south) {\FiraCodeNL ciao\Lato: linguaggio macchina};
    \node[nota, anchor=south] at ([yshift=5pt]c.south) {ciò che il programma stampa};
    \draw[wtarrow] (a.east) -- node[wtlab, above=2pt] {compilare}
      node[below=2pt, font=\FiraCodeNL\scriptsize] {gcc} (b.west);
    \draw[wtarrow] (b.east) -- node[wtlab, above=2pt] {eseguire}
      node[below=2pt, font=\FiraCodeNL\scriptsize] {./ciao} (c.west);
    \node[nota, anchor=north] at ([yshift=-4pt]a.south)
      {\textbf{l'editor}\\Visual Studio Code};
    \node[nota, anchor=north] at ([yshift=-4pt]b.south)
      {\textbf{il compilatore}\\gcc};
    \node[nota, anchor=north] at ([yshift=-4pt]c.south)
      {\textbf{il terminale}\\dove si danno i comandi};
  \end{tikzpicture}
  \caption{Dal testo scritto da noi al programma che gira: i tre strumenti
           da installare.}
  \label{fig:catena}
\end{figure}

\begin{itemize}
  \item L'\term{editor} è il quaderno su cui scriviamo il programma. Il
        programma, per ora, è un normale \textbf{file di testo} con
        estensione \code{.c} (per esempio \code{ciao.c}). Useremo
        \textbf{Visual Studio Code} (o semplicemente \textit{VS Code}), un
        editor gratuito che colora il codice, segnala gli errori mentre si
        scrive e contiene al suo interno anche un terminale.
  \item Il \term{compilatore} è il traduttore: legge il file di testo e
        produce un nuovo file, l'\term{eseguibile}, che contiene lo stesso
        programma tradotto in linguaggio macchina. Useremo \textbf{gcc}, il
        compilatore C del progetto GNU, lo stesso usato a lezione.
  \item Il \term{terminale} è una finestra in cui si \textit{scrivono} i
        comandi al computer invece di cliccare su icone. Da lì
        chiederemo a gcc di compilare e poi lanceremo il programma, che
        stamperà i suoi risultati nella stessa finestra.
\end{itemize}

Il capitolo~\ref{chap:dal-problema} spiegherà in dettaglio che cosa succede
durante la traduzione. Per ora basta l'idea: \textbf{si scrive, si compila,
si esegue}.

\section{Linux, Windows o macOS?}

Il corso \textbf{assume Linux} come sistema operativo: i comandi visti a
lezione (e all'esame) sono quelli di Linux. Non serve però cambiare
computer o cancellare Windows: in ogni sistema c'è un modo semplice per
avere gli stessi strumenti.

\begin{center}
\renewcommand{\arraystretch}{1.45}
\begin{tabular}{@{}>{\bfseries}l>{\raggedright\arraybackslash}p{7.6cm}>{\raggedright\arraybackslash}p{6.4cm}@{}}
  \toprule
  Sistema & Che cosa si installa & Perché \\
  \midrule
  Windows & \textbf{WSL} (\textit{Windows Subsystem for Linux}) con
            \textbf{Ubuntu}, e dentro Ubuntu gcc &
            WSL fa girare un vero Linux dentro Windows, senza riavvii né
            partizioni \\
  macOS   & gli \textbf{strumenti da riga di comando} di Apple
            (\textit{Command Line Tools}) &
            macOS è un parente di Linux: ha già un terminale molto
            simile, manca solo il compilatore \\
  Linux   & il pacchetto \code{build-essential} (Ubuntu, Debian, Mint) &
            è il sistema del corso: basta un comando \\
  \bottomrule
\end{tabular}
\end{center}

\textbf{VS Code} si installa allo stesso modo ovunque (passo 1). Su Windows
c'è una particolarità che conviene capire subito, perché spiega diversi
passaggi successivi: i file del corso e il compilatore stanno
\textit{dentro} Ubuntu, mentre la finestra di VS Code è un normale
programma di Windows. VS Code allora si divide in due parti che collaborano
(figura~\ref{fig:wsl}): l'interfaccia, su Windows, e un piccolo
\textit{server}, dentro Ubuntu, che legge e salva i file e lancia il
compilatore.

\begin{figure}[H]
  \centering
  \begin{tikzpicture}[font=\Lato\scriptsize,
      parte/.style={wtnode, font=\Lato\scriptsize, minimum height=1.3cm,
        text width=3.3cm}]
    % Windows
    \draw[neutral!50!white, fill=neutral!5!white, line width=0.8pt,
          rounded corners=6pt] (-8.6,2.7) rectangle (8.6,-2.6);
    \node[anchor=north west, font=\Lato\small\bfseries, text=neutral!25!black]
      at (-8.4,2.55) {%
      \tikz[baseline=0pt]{\foreach \x/\y in {0/0,0.17/0,0/0.17,0.17/0.17}
        \fill[dmaccent] (\x,\y) rectangle ++(0.15,0.15);}\ Windows};
    % VS Code
    \node[wtnode, text width=3.6cm, minimum height=3.2cm] (vsc) at (-5.9,-0.1)
      {\vscodelogo[0.8cm]\\[4pt]\textbf{Visual Studio Code}\\[2pt]
       la finestra: menu, editor, terminale};
    % Ubuntu
    \draw[orange-gradient-6!80!black, fill=orange-gradient-1!22!white,
          line width=0.8pt, rounded corners=6pt] (-2.9,1.95) rectangle (8.2,-2.25);
    \node[anchor=north west, font=\Lato\small\bfseries,
          text=orange-gradient-6!45!black] at (-2.75,1.85) {Ubuntu (dentro WSL)};
    \node[parte] (srv) at (-0.6,-0.25) {\textbf{VS Code Server}\\con l'estensione C/C++};
    \node[parte, draw=accent!60!black, fill=accent!8!white] (gcc) at (5.6,0.65) {\textbf{gcc}, \textbf{make}, \textbf{gdb}\\gli strumenti del corso};
    \node[parte, draw=green-gradient-6, fill=green-gradient-1!45!white] (file) at (5.6,-1.25) {\textbf{i tuoi file}\\{\FiraCodeNL\tiny /home/fabio/programmazione}};
    \draw[wtarrow, <->] (vsc.east) -- node[wtlab, above=2pt] {si parlano} (srv.west);
    \draw[wtarrow] (srv.east) -- node[wtlab, above, sloped] {lancia} (gcc.west);
    \draw[wtarrow] (srv.east) -- node[wtlab, below, sloped] {legge e salva} (file.west);
  \end{tikzpicture}
  \caption{Su Windows, con WSL: l'interfaccia di VS Code resta su Windows,
           mentre file e compilatore vivono in Ubuntu. Su macOS e su Linux
           tutto sta nello stesso sistema e questa divisione non c'è.}
  \label{fig:wsl}
\end{figure}

\info{Perché non un compilatore ``nativo'' per Windows?}{
  Esistono versioni di gcc per Windows (MinGW, per esempio) e funzionano.
  Ma con esse alcuni dettagli cambiano: l'eseguibile si chiama \code{a.exe}
  invece di \code{a.out}, la fine dell'input si segnala con
  \tasti{Ctrl,Z} invece di \tasti{Ctrl,D}, i comandi del terminale sono
  diversi. Con WSL si usano \textbf{esattamente} i comandi delle lezioni, ed
  è ciò che conta.}

\section{Passo 1: installare Visual Studio Code}

\subsection{Scaricare il programma}

Visual Studio Code si scarica gratuitamente dal sito ufficiale,
\urlx{https://code.visualstudio.com}. La pagina \textit{Download} propone
un pulsante per ogni sistema operativo (figura~\ref{fig:download}): si
sceglie il proprio.

\warning{Visual Studio Code, non Visual Studio}{
  Microsoft produce anche \textbf{Visual Studio} (senza ``Code''): è un
  altro programma, molto più grande, pensato per altri linguaggi. Il nome
  giusto è \textbf{Visual Studio Code}, con l'icona blu della
  figura~\ref{fig:download}.}

\begin{figure}[H]
  \centering
  \begin{minipage}{15.5cm}
  \begin{browser}[code.visualstudio.com/download]
    \centering
    {\Lato\Large\bfseries Download Visual Studio Code}\\[3pt]
    {\Lato\footnotesize\color{neutral!35!black} Free and built on open source.
     Integrated Git, debugging and extensions.}\\[8pt]
    \begin{tikzpicture}[font=\Lato\scriptsize,
        bottone/.style={fill=dmaccent, text=white, rounded corners=2pt,
          minimum width=3.3cm, minimum height=0.95cm, align=center,
          font=\Lato\scriptsize},
        dett/.style={font=\Lato\tiny, text=neutral!35!black, align=center}]
      % Windows
      \foreach \x/\y in {0/0,0.27/0,0/0.27,0.27/0.27}
        \fill[dmaccent] (-0.27+\x,0.55+\y) rectangle ++(0.25,0.25);
      \node[bottone] (w) at (0,-0.15) {{\bfseries\small Windows}\\Windows 10, 11};
      \node[dett, anchor=north] at (w.south) {User Installer \quad x64 \quad Arm64};
      % Linux
      \begin{scope}[shift={(4.6,0.85)}, scale=0.75]
        \fill[black] (0,0) ellipse (0.28 and 0.38);
        \fill[white] (0,-0.08) ellipse (0.18 and 0.26);
        \fill[white] (-0.08,0.2) circle (0.05); \fill[white] (0.08,0.2) circle (0.05);
        \fill[orange-gradient-5] (-0.07,0.12) -- (0.07,0.12) -- (0,0.03) -- cycle;
        \fill[orange-gradient-5] (-0.17,-0.37) ellipse (0.13 and 0.05);
        \fill[orange-gradient-5] (0.17,-0.37) ellipse (0.13 and 0.05);
      \end{scope}
      \node[bottone, minimum width=2.1cm] (deb) at (3.45,-0.15) {{\bfseries\small .deb}\\Debian, Ubuntu};
      \node[bottone, minimum width=2.1cm] (rpm) at (5.75,-0.15) {{\bfseries\small .rpm}\\Red Hat, Fedora};
      \node[dett, anchor=north] at ($(deb.south)!0.5!(rpm.south)$) {x64 \quad Arm64};
      % Mac
      \begin{scope}[shift={(9.6,0.68)}]
        \fill[neutral!20!black] (0,0) circle (0.17);
        \fill[neutral!20!black] (0,-0.05) ellipse (0.16 and 0.14);
        \fill[white] (0.2,0.03) circle (0.07);
        \fill[neutral!20!black, rotate around={-35:(0.03,0.24)}] (0.03,0.24) ellipse (0.035 and 0.08);
      \end{scope}
      \node[bottone] (m) at (9.6,-0.15) {{\bfseries\small Mac}\\macOS};
      \node[dett, anchor=north] at (m.south) {.zip \quad Intel \quad Apple silicon \quad Universal};
    \end{tikzpicture}
  \end{browser}
  \end{minipage}
  \caption{La pagina di download di VS Code. L'aspetto del sito cambia
           spesso; i tre gruppi di pulsanti, invece, restano.}
  \label{fig:download}
\end{figure}

\subsection{Installazione su Windows}

Il pulsante \textit{Windows} scarica un file chiamato più o meno
\code{VSCodeUserSetup-x64-1.xx.x.exe} (i numeri sono la versione). È un
normale programma di installazione:

\begin{passi}
  \item Apri il file scaricato (doppio clic, oppure clic sulla notifica di
        download del browser).
  \item Nella prima schermata scegli \textit{I accept the agreement} (accetti
        la licenza) e fai clic su \textit{Next}.
  \item Nelle due schermate successive (cartella di installazione e menu
        Start) lascia tutto com'è e fai clic su \textit{Next}.
  \item Nella schermata \textit{Select Additional Tasks}
        (figura~\ref{fig:installer}) spunta le due voci \textit{Add ``Open
        with Code'' action\dots}: aggiungono la voce ``Apri con Code'' al
        menu del tasto destro sui file e sulle cartelle, comodissima. La voce
        \textbf{Add to PATH} deve restare spuntata: permette di aprire VS
        Code scrivendo \code{code} nel terminale, e la useremo.
  \item Fai clic su \textit{Next}, poi su \textit{Install}. Al termine,
        lascia spuntato \textit{Launch Visual Studio Code} e fai clic su
        \textit{Finish}.
\end{passi}

\begin{figure}[H]
  \centering
  \begin{tikzpicture}[font=\Lato\scriptsize]
    \def\W{12.4} \def\H{7.6}
    \draw[neutral!45!white, fill=white, line width=0.6pt, rounded corners=4pt]
      (0,0) rectangle (\W,-\H);
    % barra del titolo
    \node[anchor=west] at (0.15,-0.3) {\vscodelogo[0.3cm]};
    \node[anchor=west] at (0.55,-0.3) {Setup - Microsoft Visual Studio Code (User)};
    \node[anchor=east, font=\Lato\footnotesize] at (\W-0.2,-0.3)
      {--\hspace{16pt}$\square$\hspace{16pt}\texttimes};
    % intestazione
    \node[anchor=north west, font=\Lato\footnotesize\bfseries] at (0.45,-0.72) {Select Additional Tasks};
    \node[anchor=north west] at (0.75,-1.1) {Which additional tasks should be performed?};
    \node at (\W-0.85,-1.12) {\vscodelogo[0.75cm]};
    \draw[neutral!35!white] (0,-1.68) -- (\W,-1.68);
    % corpo
    \fill[neutral!3!white] (0.02,-1.69) rectangle (\W-0.02,-6.58);
    \node[anchor=north west, text width=11cm] at (0.45,-1.9)
      {Select the additional tasks you would like Setup to perform while
       installing Visual Studio Code, then click Next.};
    \node[anchor=west] at (0.45,-2.75) {Additional icons:};
    \node[anchor=west] at (0.65,-3.12) {\wtcasellina{0}\hspace{5pt}Create a desktop icon};
    \node[anchor=west] at (0.45,-3.6) {Other:};
    \node[anchor=west] (r1) at (0.65,-3.98) {\wtcasellina{1}\hspace{5pt}Add
      \textquotedbl Open with Code\textquotedbl\ action to Windows Explorer file context menu};
    \node[anchor=west] (r2) at (0.65,-4.38) {\wtcasellina{1}\hspace{5pt}Add
      \textquotedbl Open with Code\textquotedbl\ action to Windows Explorer directory context menu};
    \node[anchor=west] at (0.65,-4.86) {\wtcasellina{1}\hspace{5pt}Register Code as an editor for supported file types};
    \node[anchor=west] (path) at (0.65,-5.34) {\wtcasellina{1}\hspace{5pt}Add to PATH (requires shell restart)};
    \draw[accent!70!black, line width=0.9pt, rounded corners=2pt]
      ([xshift=-3pt,yshift=0.5pt]r1.north west) rectangle ([xshift=3pt,yshift=-0.5pt]r2.south east);
    \draw[accent!70!black, line width=0.9pt, rounded corners=2pt]
      ([xshift=-3pt,yshift=0.5pt]path.north west) rectangle ([xshift=3pt,yshift=-0.5pt]path.south east);
    \node[anchor=west, font=\Lato\scriptsize\bfseries, text=accent!70!black, align=left]
      at ([xshift=12pt]path.east) {deve restare spuntata!};
    \node[anchor=west, font=\Lato\scriptsize\bfseries, text=accent!70!black, align=left]
      at ([xshift=8pt]r2.east |- r1.south) {comode};
    % pulsanti
    \draw[neutral!35!white] (0,-6.6) -- (\W,-6.6);
    \node[draw=neutral!50!white, fill=neutral!8!white, minimum width=1.8cm,
          minimum height=0.5cm] at (\W-5.0,-7.1) {< Back};
    \node[draw=dmaccent, line width=1pt, fill=neutral!8!white, minimum width=1.8cm,
          minimum height=0.5cm] at (\W-2.95,-7.1) {Next >};
    \node[draw=neutral!50!white, fill=neutral!8!white, minimum width=1.8cm,
          minimum height=0.5cm] at (\W-1.05,-7.1) {Cancel};
  \end{tikzpicture}
  \caption{La schermata più importante dell'installazione su Windows.}
  \label{fig:installer}
\end{figure}

\subsection{Installazione su macOS}

Il pulsante \textit{Mac} scarica un archivio \code{.zip}. Safari di solito lo
estrae da solo: nella cartella \textit{Download} compare l'applicazione
\textit{Visual Studio Code} (altrimenti basta un doppio clic sul file
\code{.zip}).

\begin{passi}
  \item Trascina \textit{Visual Studio Code} dalla cartella \textit{Download}
        alla cartella \textit{Applicazioni} (figura~\ref{fig:mac-app}): così
        è installata davvero e la si trova nel Launchpad.
  \item Aprila. Al primo avvio macOS avvisa che è un'applicazione scaricata
        da Internet: conferma con \textit{Apri}.
  \item Per poter aprire VS Code dal terminale con il comando \code{code}:
        premi \tasti{Cmd,Shift,P}, scrivi \textit{shell command} e scegli
        \textit{Shell Command: Install `code' command in PATH}.
\end{passi}

\begin{figure}[H]
  \centering
  \begin{tikzpicture}[font=\Lato\scriptsize]
    \foreach \x/\nome in {0/Download, 7.2/Applicazioni} {
      \draw[neutral!45!white, fill=white, rounded corners=5pt, line width=0.6pt]
        (\x,0) rectangle ++(4.6,-2.6);
      \fill[neutral!10!white, rounded corners=5pt] (\x,0) rectangle ++(4.6,-0.45);
      \fill[neutral!10!white] (\x,-0.25) rectangle ++(4.6,-0.2);
      \foreach \c/\k in {red-gradient-3/0, yellow-gradient-5/1, green-gradient-6/2}
        \fill[\c] (\x+0.25+\k*0.25,-0.22) circle (0.07);
      \node[font=\Lato\scriptsize\bfseries, text=neutral!30!black] at (\x+2.3,-0.22) {\nome};
    }
    \node at (2.3,-1.25) {\vscodelogo[0.9cm]};
    \node[align=center] at (2.3,-2.2) {Visual Studio Code};
    \node[text=neutral!45!black, align=center] at (9.5,-1.45)
      {(le altre applicazioni)};
    \draw[wtarrowacc, line width=1.4pt] (3.4,-1.2) to[bend left=20]
      node[wtlab, above=2pt, fill=none] {trascina} (8.0,-1.2);
  \end{tikzpicture}
  \caption{Su macOS si installa un'applicazione trascinandola nella
           cartella \textit{Applicazioni}.}
  \label{fig:mac-app}
\end{figure}

\subsection{Installazione su Linux}

Su Ubuntu, Debian e derivate si scarica il pacchetto \code{.deb} e lo si
installa dal terminale (il terminale si apre con \tasti{Ctrl,Alt,T}; il
suo uso è spiegato più avanti in questo capitolo):

\begin{console}
$ cd ~/Scaricati
$ sudo apt install ./code_*.deb
\end{console}

Il comando \code{cd} entra nella cartella dei download (su un sistema in
inglese si chiama \code{Downloads}); \code{sudo apt install} installa il
pacchetto, dopo aver chiesto la password dell'utente. In alternativa, VS
Code si trova anche nel negozio di applicazioni della distribuzione (su
Ubuntu, \textit{App Center}).

\section{Passo 2: installare il compilatore}

\subsection{Windows: WSL e Ubuntu}

Su Windows~10 (aggiornato) e Windows~11 WSL si installa con un unico
comando, da scrivere in \textbf{PowerShell} aperta \textbf{come
amministratore}, cioè con il permesso di modificare il sistema.

\begin{passi}
  \item Premi il tasto \tasto{Windows}, scrivi \textit{powershell} e, nel
        riquadro a destra, scegli \textit{Esegui come amministratore}
        (figura~\ref{fig:start}). Windows chiede conferma: rispondi
        \textit{Sì}.
  \item Nella finestra blu (o nera) che si apre, scrivi il comando
        \code{wsl -{}-install} e premi \tasto{Invio}.
  \item Attendi il download (qualche minuto: Ubuntu pesa circa un
        gigabyte) e, quando richiesto, \textbf{riavvia il computer}.
\end{passi}

\begin{figure}[H]
  \centering
  \begin{tikzpicture}[font=\Lato\scriptsize]
    \draw[neutral!40!white, fill=neutral!4!white, rounded corners=7pt,
          line width=0.6pt] (0,0) rectangle (12,-4.5);
    \draw[neutral!40!white, fill=white, rounded corners=9pt] (0.4,-0.3) rectangle (11.6,-0.88);
    \node[anchor=west] at (0.7,-0.59) {\vscico[neutral!40!black]{cerca}\hspace{6pt}powershell};
    \node[anchor=west, font=\Lato\scriptsize\bfseries] at (0.45,-1.3) {Corrispondenza migliore};
    \fill[primary!10!white, rounded corners=3pt] (0.4,-1.55) rectangle (5.6,-2.45);
    \node[fill=primary!80!black, text=white, font=\FiraCodeNL\tiny\bfseries,
          minimum size=0.55cm, rounded corners=1.5pt] at (0.95,-2.0) {>\_};
    \node[anchor=west] at (1.4,-1.82) {\textbf{Windows PowerShell}};
    \node[anchor=west, text=neutral!40!black] at (1.4,-2.15) {App};
    \fill[white, rounded corners=4pt] (6.0,-1.15) rectangle (11.6,-4.2);
    \node[fill=primary!80!black, text=white, font=\FiraCodeNL\scriptsize\bfseries,
          minimum size=0.8cm, rounded corners=2pt] at (8.8,-1.75) {>\_};
    \node[font=\Lato\footnotesize\bfseries] at (8.8,-2.45) {Windows PowerShell};
    \draw[neutral!25!white] (6.3,-2.75) -- (11.3,-2.75);
    \node[anchor=west] at (6.4,-3.0) {Apri};
    \fill[primary!10!white, rounded corners=2pt] (6.25,-3.18) rectangle (11.35,-3.6);
    \node[anchor=west, font=\Lato\scriptsize\bfseries] (adm) at (6.4,-3.39) {Esegui come amministratore};
    \node[anchor=west] at (6.4,-3.85) {Apri percorso file};
    \node at ([xshift=12pt, yshift=-2pt]adm.east) {\wtpuntatore};
  \end{tikzpicture}
  \caption{La ricerca nel menu Start: PowerShell va aperta come
           amministratore.}
  \label{fig:start}
\end{figure}

\begin{console}[Amministratore: Windows PowerShell]
PS C:\WINDOWS\system32> wsl --install
Installazione: Piattaforma macchina virtuale
Piattaforma macchina virtuale è stato installato.
Installazione: Sottosistema Windows per Linux
Sottosistema Windows per Linux è stato installato.
Installazione: Ubuntu
Ubuntu è stato installato.
L'operazione richiesta è riuscita. Le modifiche saranno effettive al riavvio del sistema.
\end{console}

\subsection{Windows: creare l'utente di Ubuntu}

I messaggi possono essere un po' diversi a seconda della versione di
Windows: quello che conta è che non compaiano errori. Dopo il riavvio si
apre da sola una finestra di Ubuntu (se non si apre, cerca \textit{Ubuntu}
nel menu Start) che completa l'installazione e chiede di creare un
\textbf{utente Linux}: un nome (in minuscolo, senza spazi; non deve essere
per forza uguale a quello di Windows) e una password.

\begin{console}[Ubuntu]
Installing, this may take a few minutes...
Please create a default UNIX user account. The username does not need to match your Windows username.
For more information visit: https://aka.ms/wslusers
Enter new UNIX username: <<fabio>>
New password:
Retype new password:
passwd: password updated successfully
Installation successful!
fabio@PC-FABIO:~$
\end{console}

\warning{Mentre scrivi la password non compare nulla}{
  Nel terminale, mentre si digita una password, \textbf{non compaiono
  nemmeno gli asterischi}: sembra che la tastiera non funzioni, invece i
  caratteri vengono letti. Scrivi la password e premi \tasto{Invio}.
  Annotala: servirà ogni volta che installerai qualcosa.}

\subsection{Windows: installare il compilatore in Ubuntu}

L'ultima riga stampata da Ubuntu, \code{fabio@PC-FABIO:\textasciitilde\$}, è il
\term{prompt}: Ubuntu è pronta a ricevere comandi (la spiegheremo pezzo per
pezzo più avanti). Restano da installare il compilatore e gli altri
strumenti del corso. Scrivi questi due comandi, uno alla volta, premendo
\tasto{Invio} dopo ciascuno:

\begin{console}[Ubuntu]
$ sudo apt update
[sudo] password for fabio:
...
$ sudo apt install build-essential gdb
...
Do you want to continue? [Y/n] <<y>>
...
\end{console}

\begin{itemize}
  \item \code{sudo} (\textit{superuser do}) esegue il comando che segue
        con i permessi di amministratore, e per questo chiede la password
        appena creata.
  \item \code{apt} è il \term{gestore di pacchetti} di Ubuntu: una specie
        di negozio di applicazioni che si usa dal terminale. Con
        \code{update} scarica l'elenco aggiornato dei programmi, con
        \code{install} installa quelli indicati (e chiede conferma: si
        risponde \code{y}, \textit{yes}).
  \item \code{build-essential} è un pacchetto che ne raccoglie altri:
        il compilatore \textbf{gcc}, il programma \textbf{make} (che useremo
        dal capitolo~\ref{chap:make}) e le \textit{librerie} standard del C.
        \textbf{gdb} è il \textit{debugger}, utile per eseguire un programma
        un'istruzione alla volta.
\end{itemize}

Per verificare che tutto sia andato a buon fine, si chiede a gcc la sua
versione:

\begin{console}[Ubuntu]
$ gcc --version
gcc (Ubuntu 13.3.0-6ubuntu2~24.04.1) 13.3.0
Copyright (C) 2023 Free Software Foundation, Inc.
This is free software; see the source for copying conditions.  There is NO
warranty; not even for MERCHANTABILITY or FITNESS FOR A PARTICULAR PURPOSE.
\end{console}

I numeri di versione possono essere diversi: va bene qualunque versione.
Se invece compare \code{gcc: command not found}, l'installazione non è
riuscita: ripeti i due comandi \code{apt} e controlla che non stampino
errori.

\warning{Se \code{wsl -{}-install} non funziona}{
  L'errore più comune all'avvio di Ubuntu è un messaggio che parla di
  \textit{virtualizzazione} (per esempio il codice \code{0x80370102}): WSL
  usa una funzione del processore che su alcuni computer è disattivata.
  Si attiva dal BIOS/UEFI del computer (la voce si chiama
  \textit{Virtualization Technology}, \textit{Intel VT-x} o \textit{AMD-V});
  la procedura cambia da modello a modello, quindi conviene cercare
  ``attivare virtualizzazione'' seguito dal modello del proprio PC. Le
  istruzioni ufficiali di Microsoft sono indicate nei riferimenti a fine
  capitolo.}

\subsection{macOS: gli strumenti da riga di comando}

Su macOS il compilatore fa parte degli \textit{strumenti da riga di
comando} di Apple, gratuiti. Non serve installare Xcode (l'ambiente di
sviluppo completo, molto più grande): bastano gli strumenti.

\begin{passi}
  \item Apri il \textbf{Terminale}: premi \tasti{Cmd,Spazio} (si apre
        Spotlight), scrivi \textit{terminale} e premi \tasto{Invio}.
  \item Scrivi il comando \code{xcode-select -{}-install} e premi
        \tasto{Invio}.
  \item Compare una finestra come quella della figura~\ref{fig:xcode}: fai
        clic su \textit{Installa} (non su \textit{Ottieni Xcode}) e accetta
        la licenza. Il download richiede qualche minuto.
\end{passi}

\begin{figure}[H]
  \centering
  \begin{tikzpicture}[font=\Lato\scriptsize]
    \draw[neutral!35!white, fill=neutral!3!white, rounded corners=9pt,
          line width=0.6pt] (0,0) rectangle (10.2,-3.3);
    \node[fill=neutral!25!white, rounded corners=4pt, minimum size=1cm,
          font=\FiraCodeNL\small\bfseries, text=neutral!25!black] at (0.95,-1.05) {>\_};
    \node[anchor=north west, text width=8.3cm, align=left] at (1.65,-0.35)
      {\textbf{Il comando ``xcode-select'' richiede gli strumenti per
       sviluppatori della riga di comando. Vuoi installarli adesso?}\\[4pt]
       Scegli Installa per scaricare e installare ora gli strumenti per
       sviluppatori della riga di comando.};
    \node[draw=neutral!45!white, fill=white, rounded corners=3pt,
          minimum width=2.1cm, minimum height=0.48cm] at (3.0,-2.75) {Ottieni Xcode};
    \node[draw=neutral!45!white, fill=white, rounded corners=3pt,
          minimum width=1.8cm, minimum height=0.48cm] at (7.0,-2.75) {Non ora};
    \node[fill=dmaccent, text=white, rounded corners=3pt,
          minimum width=1.8cm, minimum height=0.48cm] (ok) at (9.0,-2.75) {Installa};
    \node at ([xshift=6pt, yshift=-5pt]ok.center) {\wtpuntatore};
  \end{tikzpicture}
  \caption{La richiesta di installazione degli strumenti da riga di
           comando (il testo esatto dipende dalla versione di macOS).}
  \label{fig:xcode}
\end{figure}

Al termine, la verifica:

\begin{console}[Terminale di macOS]
$ gcc --version
Apple clang version 16.0.0 (clang-1600.0.26.6)
Target: arm64-apple-darwin24.1.0
Thread model: posix
InstalledDir: /Library/Developer/CommandLineTools/usr/bin
\end{console}

Sorpresa: il comando \code{gcc} risponde \textit{Apple clang}. Su macOS il
compilatore è \textbf{clang}, un altro compilatore C, che per comodità
risponde anche al nome \code{gcc} e accetta le stesse opzioni. Per tutto
ciò che si fa nel corso il comportamento è lo stesso.

\info{Il prompt di macOS}{
  Nel Terminale di macOS il prompt finisce con il simbolo \code{\%} invece
  che con \code{\$} (per esempio \code{fabio@MacBook-Air \textasciitilde\
  \%}). In questa dispensa i comandi sono sempre preceduti da \code{\$}: su
  Mac si scrivono allo stesso modo.}

\subsection{Linux}

Su Ubuntu, Debian e Linux Mint bastano gli stessi due comandi visti per
Windows (in Ubuntu, il terminale si apre con \tasti{Ctrl,Alt,T}):

\begin{console}
$ sudo apt update
$ sudo apt install build-essential gdb
\end{console}

Su Fedora il comando equivalente è
\code{sudo dnf install gcc make gdb}. La verifica, anche qui, è
\code{gcc -{}-version}.

\section{Passo 3: le estensioni di VS Code}

VS Code, appena installato, sa già colorare il codice C. Le
\term{estensioni} sono componenti aggiuntivi gratuiti che gli insegnano a
fare di più. Si installano dalla vista \textit{Extensions}: si apre con
l'icona dei quattro quadratini nella barra a sinistra (oppure con
\tasti{Ctrl,Shift,X}, su Mac \tasti{Cmd,Shift,X}); si scrive il nome
nella casella di ricerca e si fa clic su \pulsante{Install}
(figura~\ref{fig:estensioni}).

\begin{figure}[H]
  \centering
  \begin{tikzpicture}[font=\Lato\scriptsize,
      esten/.style={regular polygon, regular polygon sides=6, minimum size=0.95cm,
        inner sep=0pt, font=\Lato\tiny\bfseries, text=white},
      nome/.style={anchor=west, text=dmfg, font=\Lato\scriptsize\bfseries},
      desc/.style={anchor=west, text=dmdim, font=\Lato\tiny},
      edit/.style={anchor=west, text=dmfg, font=\Lato\tiny}]
    \def\W{10.2}
    \fill[dmbar] (0,0) rectangle (\W,-5.4);
    \draw[dmborder] (0.75,0) -- (0.75,-5.4);
    % barra delle attivita'
    \foreach \ico/\y in {esplora/-0.55, cerca/-1.2, git/-1.85, debug/-2.5} {
      \node at (0.375,\y) {\vscico{\ico}};
    }
    \node at (0.375,-3.15) {\vscico[white]{estensioni}};
    \fill[white] (0,-2.92) rectangle (0.04,-3.38);
    % intestazione e ricerca
    \node[anchor=west, text=dmfg, font=\Lato\tiny] at (0.95,-0.3) {EXTENSIONS: MARKETPLACE};
    \draw[dmaccent, fill=dminput, line width=0.6pt] (0.95,-0.58) rectangle (\W-0.25,-1.02);
    \node[anchor=west, text=dmfg] at (1.05,-0.8) {c/c++};
    % risultati
    \fill[dmsel] (0.76,-1.2) rectangle (\W,-2.5);
    \foreach \y/\c in {-1.2/purple-gradient-6!80!black, -2.5/purple-gradient-5!70!black,
                       -3.8/purple-gradient-4!60!black}
      \node[esten, fill=\c] at (1.5,\y-0.62) {C/C++};
    \node[nome] at (2.2,-1.45) {C/C++};
    \node[desc] at (2.2,-1.78) {C/C++ IntelliSense, debugging, and code browsing.};
    \node[edit] (p1) at (2.2,-2.13) {\tikz[baseline=-0.6ex]\node[circle, fill=dmaccent, inner sep=0.6pt,
      text=white, font=\Lato\tiny]{\itick};\ Microsoft};
    \node[nome] at (2.2,-2.75) {C/C++ Extension Pack};
    \node[desc] at (2.2,-3.08) {Popular extensions for C++ development in Visual Studio Code.};
    \node[edit] at (2.2,-3.43) {\tikz[baseline=-0.6ex]\node[circle, fill=dmaccent, inner sep=0.6pt,
      text=white, font=\Lato\tiny]{\itick};\ Microsoft};
    \node[nome] at (2.2,-4.05) {C/C++ Themes};
    \node[desc] at (2.2,-4.38) {UI Themes for C/C++ extension.};
    \node[edit] at (2.2,-4.73) {\tikz[baseline=-0.6ex]\node[circle, fill=dmaccent, inner sep=0.6pt,
      text=white, font=\Lato\tiny]{\itick};\ Microsoft};
    \foreach \y in {-2.13,-3.43,-4.73}
      \node[fill=dmaccent, text=white, font=\Lato\tiny, inner xsep=4pt, inner ysep=2pt,
            rounded corners=1pt] at (\W-0.75,\y) {Install};
    \node at (\W-0.45,-2.25) {\wtpuntatore};
    % didascalia a destra
    \node[anchor=west, font=\Lato\scriptsize, text=accent!70!black, align=left]
      at (\W+0.3,-1.85) {\textbf{questa}: l'estensione\\ufficiale di Microsoft};
    \draw[wtarrowacc] (\W+0.3,-1.85) -- (\W+0.05,-1.85);
  \end{tikzpicture}
  \caption{La vista \textit{Extensions} con la ricerca ``c/c++''. Basta la
           prima: \textit{C/C++} di Microsoft (si riconosce dal segno blu di
           ``editore verificato'').}
  \label{fig:estensioni}
\end{figure}

Le estensioni da installare sono poche:

\begin{center}
\renewcommand{\arraystretch}{1.45}
\begin{tabular}{@{}>{\raggedright\arraybackslash}p{3.2cm}>{\raggedright\arraybackslash}p{2.4cm}>{\raggedright\arraybackslash}p{10.6cm}@{}}
  \toprule
  \textbf{Estensione} & \textbf{Per chi} & \textbf{A che cosa serve} \\
  \midrule
  \textbf{C/C++} (Microsoft) & tutti &
    \textit{IntelliSense}: suggerimenti mentre si scrive (nomi di funzione,
    parametri), errori sottolineati in rosso ancora prima di compilare,
    navigazione nel codice (clic destro, \textit{Go to Definition}),
    debugger. \\
  \textbf{WSL} (Microsoft) & solo Windows &
    collega VS Code a Ubuntu, come nella figura~\ref{fig:wsl}. \\
  \bottomrule
\end{tabular}
\end{center}

Spesso non serve nemmeno cercarle: aprendo per la prima volta un file
\code{.c}, VS Code propone da solo, con una notifica in basso a destra, di
installare l'estensione consigliata per il linguaggio C.

\info{E la lingua italiana?}{
  VS Code può essere tradotto in italiano con l'estensione
  \textit{Italian Language Pack}. In questa dispensa però i menu sono
  citati \textbf{in inglese}, come compaiono dopo l'installazione: se
  scegli l'italiano dovrai tradurre i nomi delle voci (\textit{File} resta
  \textit{File}, \textit{Terminal} diventa \textit{Terminale}, e così
  via).}

\subsection{Windows: collegare VS Code a Ubuntu}

Su Windows l'ordine conta: prima si installa l'estensione \textbf{WSL}, poi
si apre VS Code \textit{dentro} Ubuntu, e solo allora si installa
\textbf{C/C++}. Il modo più semplice per aprire VS Code dentro Ubuntu è
farlo dal terminale di Ubuntu, con il comando \code{code} seguito da un
punto (il punto significa ``la cartella in cui mi trovo''):

\begin{console}[Ubuntu]
$ code .
Installing VS Code Server for Linux x64 (...)
Downloading: 100%
Unpacking: 100%
\end{console}

La prima volta VS Code installa il suo \textit{server} dentro Ubuntu (lo
fa da solo, in pochi secondi), poi si apre la finestra. Come si capisce che
è collegata a Ubuntu? Dall'\term{indicatore remoto}, il riquadro blu
nell'angolo in basso a sinistra (figura~\ref{fig:remoto}).

\begin{figure}[H]
  \centering
  \begin{tikzpicture}[font=\Lato\scriptsize]
    \fill[dmbar] (0,0) rectangle (13,-0.44);
    \draw[dmborder] (0,0) -- (13,0);
    \fill[dmaccent] (0,0) rectangle (2.5,-0.44);
    \node[anchor=west, text=white] (rem) at (0.15,-0.22)
      {\vscico[white]{remoto}\hspace{5pt}WSL: Ubuntu};
    \node[anchor=west, text=dmfg] at (2.75,-0.22) {\icoerr\ 0\hspace{6pt}\icowarn\ 0};
    \node[anchor=east, text=dmfg] at (12.8,-0.22)
      {Ln 5, Col 31\hspace{14pt}Spaces: 4\hspace{14pt}UTF-8\hspace{14pt}LF\hspace{14pt}\{\} C\hspace{14pt}Linux};
    \draw[wtarrowacc] (1.25,-1.15) -- (1.25,-0.5);
    \node[anchor=north, font=\Lato\scriptsize, text=accent!70!black, align=center]
      at (1.25,-1.15) {\textbf{collegato a Ubuntu}};
  \end{tikzpicture}
  \caption{La barra di stato di VS Code collegato a Ubuntu. Se il riquadro
           in basso a sinistra non dice \textit{WSL: Ubuntu}, VS Code sta
           lavorando su Windows e non troverà il compilatore.}
  \label{fig:remoto}
\end{figure}

Ora, nella vista \textit{Extensions}, l'estensione C/C++ mostra il
pulsante \pulsante{Install in WSL: Ubuntu}: è quello giusto, perché
l'estensione deve lavorare dove stanno i file e il compilatore. Le volte
successive, VS Code ricorda le cartelle aperte: dal menu \menu{File,Open
Recent} si riapre direttamente una cartella di Ubuntu, già collegata.

\info{Dove sono i file di Ubuntu?}{
  La cartella personale di Ubuntu (\code{/home/fabio}) è visibile anche da
  Esplora file di Windows, alla voce \textit{Linux} nel riquadro di
  sinistra. Viceversa, da Ubuntu il disco \code{C:} di Windows si trova in
  \code{/mnt/c}. Conviene però tenere i file del corso \textbf{dentro
  Ubuntu}, nella cartella personale: lì compilare è più veloce e non ci
  sono problemi di permessi.}

\section{Passo 4: la cartella del corso}

Conviene tenere tutti i programmi del corso in una cartella dedicata. La
creiamo dal terminale (su Windows, quello di \textbf{Ubuntu}) e la apriamo
in VS Code:

\begin{console}
$ mkdir programmazione
$ cd programmazione
$ code .
\end{console}

\code{mkdir} (\textit{make directory}) crea la cartella, \code{cd}
(\textit{change directory}) ci entra, \code{code .} apre VS Code sulla
cartella corrente. In alternativa si può fare tutto con il mouse: si crea
la cartella come al solito e la si apre in VS Code da \menu{File,Open
Folder...}.

La prima volta che si apre una cartella, VS Code chiede se ci si fida
degli autori dei file che contiene (\textit{Do you trust the authors of the
files in this folder?}): sono i nostri file, quindi si risponde
\pulsante{Yes, I trust the authors}. Senza questa conferma VS Code apre la
cartella in \textit{modalità con restrizioni} e disattiva le estensioni.

\subsection{Com'è fatta la finestra di VS Code}

La figura~\ref{fig:vscode} mostra VS Code con la cartella del corso
aperta, un file C nell'editor e il terminale in basso. Vale la pena
guardarla con calma: sono le zone che useremo sempre.

% L'editor della figura contiene un listato vero (senza barra delle schede)
\newsavebox{\vsceditor}
\begin{lrbox}{\vsceditor}
\begin{minipage}{12.6cm}
\lstset{backgroundcolor=\color{dmbg}, rulecolor=\color{dmbg}}
\begin{lstlisting}[style=c, aboveskip=0pt, belowskip=0pt]
#include <stdio.h>

int main(void)
{
    printf("Ciao, mondo!\n");
    return 0;
}
\end{lstlisting}
\end{minipage}
\end{lrbox}

\begin{figure}[H]
  \centering
  \begin{tikzpicture}[font=\Lato\scriptsize,
      num/.style={circle, fill=accent!75!black, text=white, draw=white,
        line width=0.8pt, font=\Lato\scriptsize\bfseries, inner sep=0pt,
        minimum size=0.42cm},
      piccolo/.style={font=\Lato\tiny, text=dmfg, inner sep=0pt},
      term/.style={anchor=west, font=\FiraCodeNL\scriptsize, inner sep=0pt}]
    \def\W{17.4} \def\H{10.3}
    % fondo e barra del titolo
    \fill[dmbar, rounded corners=4pt] (0,0) rectangle (\W,-\H);
    \node at (0.38,-0.25) {\vscodelogo[0.3cm]};
    \node[piccolo, anchor=west] at (0.8,-0.25)
      {File\hspace{9pt}Edit\hspace{9pt}Selection\hspace{9pt}View\hspace{9pt}Go\hspace{9pt}Run\hspace{9pt}Terminal\hspace{9pt}Help};
    \fill[dminput, rounded corners=3pt] (6.6,-0.1) rectangle (11.0,-0.41);
    \node[piccolo, text=dmdim] at (8.8,-0.255) {\scalebox{0.7}{\vscico{cerca}}\ programmazione [WSL: Ubuntu]};
    \node[piccolo, anchor=east, font=\Lato\scriptsize] at (\W-0.25,-0.25)
      {--\hspace{15pt}$\square$\hspace{15pt}\texttimes};
    \draw[dmborder] (0,-0.5) -- (\W,-0.5);
    % barra delle attivita'
    \draw[dmborder] (0.75,-0.5) -- (0.75,-\H+0.35);
    \node at (0.375,-0.95) {\vscico[white]{esplora}};
    \fill[dmaccent] (0,-0.72) rectangle (0.045,-1.18);
    \foreach \ico/\y in {cerca/-1.6, git/-2.25, debug/-2.9, estensioni/-3.55}
      \node at (0.375,\y) {\vscico{\ico}};
    \node at (0.375,-\H+1.4) {\vscico{account}};
    \node at (0.375,-\H+0.8) {\vscico{impostazioni}};
    % barra laterale (Explorer)
    \draw[dmborder] (4.1,-0.5) -- (4.1,-\H+0.35);
    \node[piccolo, anchor=west] at (0.95,-0.78) {EXPLORER};
    \node[piccolo, anchor=east] at (3.95,-0.78) {$\cdots$};
    \node[piccolo, anchor=west, font=\Lato\tiny\bfseries] at (0.9,-1.15)
      {\rotatebox[origin=c]{-90}{›}\ PROGRAMMAZIONE [WSL: UBUNTU]};
    \fill[dmsel] (0.76,-1.38) rectangle (4.09,-1.72);
    \draw[dmaccent, line width=0.5pt] (0.77,-1.38) rectangle (4.08,-1.72);
    \node[piccolo, anchor=west, font=\Lato\scriptsize] at (1.15,-1.55)
      {{\FiraCodeNL\bfseries\color{vsciconblue}C}\ \ ciao.c};
    \node[piccolo, anchor=west, font=\Lato\scriptsize] at (1.15,-1.92)
      {{\color{vsctabdim}\ifilecode}\ \ ciao};
    % editor: barra delle schede
    \fill[dmbg] (4.11,-0.51) rectangle (\W,-\H+3.75);
    \fill[dmbar] (4.11,-0.51) rectangle (\W,-0.98);
    \fill[dmbg] (4.11,-0.51) rectangle (6.3,-0.98);
    \fill[dmaccent] (4.11,-0.51) rectangle (6.3,-0.55);
    \node[piccolo, anchor=west, font=\Lato\scriptsize] at (4.3,-0.75)
      {{\FiraCodeNL\bfseries\color{vsciconblue}C}\ \ ciao.c\hspace{9pt}{\color{dmdim}\texttimes}};
    \node at (\W-1.05,-0.75) {\vscico[dmfg]{gioca}};
    \node[piccolo] at (\W-0.7,-0.75) {\rotatebox[origin=c]{-90}{›}};
    \node[piccolo, anchor=west, text=dmdim] at (4.3,-1.2)
      {{\FiraCodeNL\color{vsciconblue}C}\ ciao.c\ \ ›\ \ main};
    \node[anchor=north west, inner sep=0pt] at (4.15,-1.45) {\usebox{\vsceditor}};
    % pannello del terminale
    \draw[dmborder] (4.1,-\H+3.75) -- (\W,-\H+3.75);
    \fill[dmbar] (4.11,-\H+3.74) rectangle (\W,-\H+0.36);
    \node[piccolo, anchor=west, text=dmdim] (t1) at (4.35,-\H+3.42) {PROBLEMS};
    \node[piccolo, anchor=west, text=dmdim] (t2) at ([xshift=12pt]t1.east) {OUTPUT};
    \node[piccolo, anchor=west, text=dmdim] (t3) at ([xshift=12pt]t2.east) {DEBUG CONSOLE};
    \node[piccolo, anchor=west, text=white] (t4) at ([xshift=12pt]t3.east) {TERMINAL};
    \node[piccolo, anchor=west, text=dmdim] at ([xshift=12pt]t4.east) {PORTS};
    \fill[dmaccent] ([yshift=-3pt]t4.south west) rectangle ([yshift=-3.6pt]t4.south east);
    \node[piccolo, anchor=east] at (\W-0.3,-\H+3.42)
      {\iterminal\hspace{3pt}bash\hspace{10pt}+\hspace{10pt}\texttimes};
    \node[term] at (4.35,-\H+2.85) {\textcolor{termgreen}{\bfseries fabio@PC-FABIO}\textcolor{dmfg}{:}\textcolor{termblue}{\bfseries\textasciitilde/programmazione}\textcolor{dmfg}{\$ }\textcolor{white}{gcc -Wall -o ciao ciao.c}};
    \node[term] at (4.35,-\H+2.45) {\textcolor{termgreen}{\bfseries fabio@PC-FABIO}\textcolor{dmfg}{:}\textcolor{termblue}{\bfseries\textasciitilde/programmazione}\textcolor{dmfg}{\$ }\textcolor{white}{./ciao}};
    \node[term, text=dmfg] at (4.35,-\H+2.05) {Ciao, mondo!};
    \node[term] at (4.35,-\H+1.65) {\textcolor{termgreen}{\bfseries fabio@PC-FABIO}\textcolor{dmfg}{:}\textcolor{termblue}{\bfseries\textasciitilde/programmazione}\textcolor{dmfg}{\$ }\textcolor{white}{\rule[-1pt]{4pt}{7pt}}};
    % barra di stato
    \draw[dmborder] (0,-\H+0.35) -- (\W,-\H+0.35);
    \begin{scope}
      \clip[rounded corners=4pt] (0,0) rectangle (\W,-\H);
      \fill[dmaccent] (0,-\H+0.35) rectangle (2.3,-\H);
    \end{scope}
    \node[piccolo, anchor=west, text=white, font=\Lato\tiny] at (0.12,-\H+0.175)
      {\scalebox{0.6}{\vscico[white]{remoto}}\ \ WSL: Ubuntu};
    \node[piccolo, anchor=west] at (2.5,-\H+0.175) {\icoerr\ 0\hspace{5pt}\icowarn\ 0};
    \node[piccolo, anchor=east] at (\W-0.25,-\H+0.175)
      {Ln 5, Col 31\hspace{10pt}Spaces: 4\hspace{10pt}UTF-8\hspace{10pt}LF\hspace{10pt}\{\} C\hspace{10pt}Linux};
    % numeri
    \node[num] at (0.375,-4.4) {1};
    \node[num] at (2.45,-3.0) {2};
    \node[num] at (13.6,-3.0) {3};
    \node[num] at (\W-1.75,-0.75) {4};
    \node[num] at (14.8,-\H+2.2) {5};
    \node[num] at (2.6,-\H+0.6) {6};
  \end{tikzpicture}
  \caption{Visual Studio Code con la cartella \code{programmazione} aperta
           (su Windows, collegato a Ubuntu).}
  \label{fig:vscode}
\end{figure}

\begin{passi}
  \item La \term{barra delle attività}: cambia ciò che compare nella barra
        laterale. Le icone, dall'alto: \textit{Explorer} (i file),
        \textit{Search} (cerca nei file), \textit{Source Control} (Git),
        \textit{Run and Debug} (il debugger), \textit{Extensions}.
  \item La \term{barra laterale}, qui con l'\textit{Explorer}: i file della
        cartella aperta. Un clic apre un file nell'editor; il pulsante con
        il foglio e il segno ``+'', che compare passando con il mouse sul
        nome della cartella, crea un nuovo file.
  \item L'\term{editor}: il file aperto, con i numeri di riga a sinistra. In
        alto ci sono le schede dei file aperti e il ``percorso'' (\textit{ciao.c
        › main}: il cursore è dentro la funzione \code{main}).
  \item Il pulsante \textit{Run}, aggiunto dall'estensione C/C++ (ne
        parleremo alla fine del capitolo).
  \item Il \term{pannello} in basso, qui con il \textit{terminale}
        integrato: un vero terminale (su Windows, di Ubuntu) già posizionato
        nella cartella aperta. Si apre dal menu \menu{Terminal,New
        Terminal}; accanto alla voce del menu è indicata anche la
        scorciatoia da tastiera, che dipende dalla tastiera in uso.
  \item La \term{barra di stato}: a sinistra l'indicatore remoto e il
        numero di errori e avvisi; a destra la posizione del cursore
        (\textit{Ln} riga, \textit{Col} colonna) e il linguaggio
        riconosciuto (\textit{C}).
\end{passi}

\subsection{Creare un file}

Un nuovo file si crea dall'\textit{Explorer} (pulsante ``nuovo file''
accanto al nome della cartella) oppure da \menu{File,New File...}. Il nome
deve finire con \textbf{\code{.c}}: è l'\term{estensione} del nome a dire a
VS Code (e a gcc) che il file contiene codice C. Con un nome come
\code{ciao.txt}, o senza estensione, niente colori, niente suggerimenti e
gcc si rifiuterebbe di compilarlo.

\warning{Nomi di file semplici}{
  Usa nomi brevi, in minuscolo, \textbf{senza spazi} e senza lettere
  accentate: \code{esercizio1.c}, \code{media\_voti.c}. Gli spazi nei nomi
  costringono a usare le virgolette nei comandi del terminale e sono una
  fonte continua di errori.}

\section{Il terminale}

Il terminale è un dialogo scritto con il computer: noi scriviamo un
\textbf{comando} e premiamo \tasto{Invio}, il computer lo esegue e risponde
con del testo. Il programma che legge i comandi si chiama \term{shell}; in
Ubuntu è \textit{bash}, su macOS \textit{zsh}. Per il corso servono pochi
comandi, ma vanno capiti bene: sono il modo in cui compileremo ed eseguiremo
ogni programma.

\subsection{Il prompt}

Quando la shell è pronta mostra il \term{prompt}, una riga che finisce con
\code{\$} (su Mac \code{\%}) e dopo la quale si scrive. Il prompt di
Ubuntu non è decorativo: dice \textit{chi} siamo, \textit{dove} siamo e
che la shell \textit{aspetta} un comando (figura~\ref{fig:prompt}).

\begin{figure}[H]
  \centering
  \begin{tikzpicture}[font=\FiraCodeNL\normalsize,
      pezzo/.style={inner sep=0pt, anchor=base west},
      eti/.style={font=\Lato\scriptsize, text=neutral!20!black, align=center}]
    \fill[consolebg, rounded corners=4pt] (-0.4,0.55) rectangle (8.6,-0.4);
    \node[pezzo, text=termgreen, font=\FiraCodeNL\normalsize\bfseries] (u) at (0,0) {fabio};
    \node[pezzo, text=termgreen, font=\FiraCodeNL\normalsize\bfseries] (at) at (u.base east) {@};
    \node[pezzo, text=termgreen, font=\FiraCodeNL\normalsize\bfseries] (h) at (at.base east) {PC-FABIO};
    \node[pezzo, text=consolefg] (due) at (h.base east) {:};
    \node[pezzo, text=termblue, font=\FiraCodeNL\normalsize\bfseries] (d) at (due.base east) {\textasciitilde/programmazione};
    \node[pezzo, text=consolefg] (s) at (d.base east) {\$};
    \node[pezzo, text=white, font=\FiraCodeNL\normalsize\bfseries] (c) at ([xshift=8pt]s.base east) {ls -l};
    % graffe sotto (punta in basso) per le parti del prompt
    \foreach \p in {u, h, d, s}
      \draw[wtbrace] ([yshift=-0.55cm]\p.base west) -- ([yshift=-0.55cm]\p.base east);
    \node[eti, anchor=north] at ([yshift=-0.75cm]u.base) {l'utente};
    \node[eti, anchor=north] at ([yshift=-0.75cm]h.base) {il nome del\\computer};
    \node[eti, anchor=north] at ([yshift=-0.75cm]d.base) {la \textbf{cartella corrente}\\(\textasciitilde\ è la tua cartella\\personale)};
    \node[eti, anchor=north] at ([yshift=-0.75cm]s.base) {``pronto''};
    % graffa sopra (punta in alto) per il comando
    \draw[wtbrace] ([yshift=0.75cm]c.base east) -- ([yshift=0.75cm]c.base west);
    \node[eti, anchor=south] at ([yshift=0.95cm]c.base) {il comando, scritto da te};
  \end{tikzpicture}
  \caption{Le parti del prompt di Ubuntu.}
  \label{fig:prompt}
\end{figure}

La parte più importante è la \term{cartella corrente} (o \textit{directory
di lavoro}): ogni comando lavora \textit{lì}, a meno di indicare
diversamente. Se scriviamo \code{gcc ciao.c} e il file \code{ciao.c} non
sta nella cartella corrente, gcc non lo trova. Moltissimi problemi dei
principianti si risolvono guardando il prompt.

\info{Il prompt, nella dispensa}{
  D'ora in poi i comandi sono mostrati con il solo \code{\$} come prompt e
  in \textbf{grassetto bianco}; le righe normali sono ciò che il computer
  risponde, e il testo in \textbf{\color{orange-gradient-6!80!black}giallo}
  è ciò che si scrive mentre un programma è in esecuzione. Il \code{\$}
  \textbf{non va digitato}: si scrive solo ciò che lo segue.}

\subsection{Cartelle e percorsi}

I file sono organizzati in cartelle, che contengono altre cartelle, che
contengono altri file: un \term{albero}, che in Linux parte da un'unica
radice chiamata \code{/} (figura~\ref{fig:albero}).

\begin{figure}[H]
  \centering
  \begin{tikzpicture}[font=\Lato\scriptsize]
    \node[fsdir] (root) at (0,0) {/};
    \node[fsdir] (home) at (-4.2,-1.1) {home};
    \node[fsdir] (usr) at (0.6,-1.1) {usr};
    \node[fsdir] (mnt) at (4.4,-1.1) {mnt};
    \node[fsdir] (fabio) at (-4.2,-2.2) {fabio};
    \node[fsqui] (prog) at (-4.2,-3.3) {programmazione};
    \node[fsfile] (ciao) at (-5.6,-4.4) {ciao.c};
    \node[fsfile] (eseg) at (-2.8,-4.4) {ciao};
    \node[fsdir] (bin) at (0.6,-2.2) {bin};
    \node[fsfile] (gcc) at (0.6,-3.3) {gcc};
    \node[fsdir] (c) at (4.4,-2.2) {c};
    \foreach \a/\b in {root/home, root/usr, root/mnt, home/fabio, fabio/prog,
                       prog/ciao, prog/eseg, usr/bin, bin/gcc, mnt/c}
      \draw[fsramo] (\a.south) -- ++(0,-0.2) -| (\b.north);
    \node[mnota, anchor=west, text=primary!50!black] at ([xshift=4pt]fabio.east)
      {\FiraCodeNL\textasciitilde\Lato: la tua cartella personale};
    \node[mnota, anchor=west, text=orange-gradient-6!50!black] at ([xshift=4pt]prog.east)
      {la cartella corrente};
    \node[mnota, anchor=north, align=center] at ([yshift=-2pt]gcc.south)
      {il programma gcc\\è un file anche lui};
    \node[mnota, anchor=north, align=center] at ([yshift=-2pt]c.south)
      {il disco \FiraCodeNL C:\Lato\ di Windows\\(solo con WSL)};
  \end{tikzpicture}
  \caption{Una parte dell'albero delle cartelle di Ubuntu.}
  \label{fig:albero}
\end{figure}

Un file si indica con il suo \term{percorso}, cioè la strada per
raggiungerlo, con i nomi delle cartelle separati da \code{/}:

\begin{itemize}
  \item un percorso \term{assoluto} parte dalla radice e vale ovunque ci si
        trovi: \code{/home/fabio/programmazione/ciao.c};
  \item un percorso \term{relativo} parte dalla cartella corrente: se si è
        in \code{/home/fabio}, lo stesso file è
        \code{programmazione/ciao.c}; se si è già in \code{programmazione},
        è semplicemente \code{ciao.c}.
\end{itemize}

Tre nomi speciali completano il quadro: \code{.} è la cartella corrente,
\code{..} è la cartella che la contiene (``sali di un livello'') e
\code{\textasciitilde} è la cartella personale. Dalla cartella
\code{programmazione}, quindi, \code{..} è \code{/home/fabio}.

\subsection{I comandi essenziali}

\begin{center}
\renewcommand{\arraystretch}{1.4}
\begin{tabular}{@{}>{\FiraCodeNL\small}l>{\raggedright\arraybackslash}p{6.0cm}>{\raggedright\arraybackslash}p{5.6cm}@{}}
  \toprule
  \normalfont\textbf{Comando} & \textbf{Che cosa fa} & \textbf{Da dove viene il nome} \\
  \midrule
  pwd           & stampa il percorso della cartella corrente & \textit{print working directory} \\
  ls            & elenca il contenuto della cartella corrente & \textit{list} \\
  ls -l         & come sopra, con dimensione e data dei file & \textit{long} \\
  cd nome       & entra nella cartella \code{nome} & \textit{change directory} \\
  cd ..         & sale nella cartella superiore & \\
  cd            & torna nella cartella personale & \\
  mkdir nome    & crea la cartella \code{nome} & \textit{make directory} \\
  cp da a       & copia il file \code{da} nel file \code{a} & \textit{copy} \\
  mv da a       & sposta o rinomina un file & \textit{move} \\
  rm file       & cancella un file (\textbf{senza cestino!}) & \textit{remove} \\
  cat file      & mostra il contenuto di un file di testo & \textit{concatenate} \\
  clear         & pulisce lo schermo del terminale & \\
  \bottomrule
\end{tabular}
\end{center}

Ecco una breve sessione di lavoro; la spiegazione segue subito dopo.

\begin{console}
$ pwd
/home/fabio
$ ls
programmazione
$ cd programmazione
$ ls
ciao.c
$ mkdir lezione2
$ ls
ciao.c  lezione2
$ cp ciao.c lezione2/prova.c
$ cd lezione2
$ ls
prova.c
$ cd ..
$ pwd
/home/fabio/programmazione
\end{console}

Si parte dalla cartella personale (\code{pwd} lo conferma), si entra in
\code{programmazione}, si crea la sottocartella \code{lezione2} e ci si
copia il file \code{ciao.c} con il nuovo nome \code{prova.c} (si noti il
percorso relativo \code{lezione2/prova.c}). Poi si entra in
\code{lezione2}, si controlla che la copia ci sia, e con \code{cd ..} si
torna su di un livello. Molti comandi, quando riescono, \textbf{non
stampano niente}: è la convenzione di Linux, ``nessuna notizia, buona
notizia''.

\error{\code{rm} non perdona}{
  Il comando \code{rm} cancella definitivamente: il file non passa per il
  cestino e non si recupera. Controlla sempre due volte il nome prima di
  premere \tasto{Invio}.}

\subsection{Trucchi che fanno risparmiare tempo}

\begin{itemize}
  \item \tasto{Tab} \textbf{completa} i nomi: scrivendo \code{cd prog} e
        premendo \tasto{Tab}, la shell completa da sola in
        \code{cd programmazione/}. Oltre a risparmiare tempo evita gli
        errori di battitura: se il completamento non funziona, il file non
        si trova dove pensiamo.
  \item \tasto{$\uparrow$} e \tasto{$\downarrow$} scorrono i \textbf{comandi già dati}: per
        ricompilare dopo una modifica basta premere \tasto{$\uparrow$} fino a
        ritrovare il comando \code{gcc} e poi \tasto{Invio}.
  \item \tasti{Ctrl,C} \textbf{interrompe} il programma in esecuzione. Se un
        programma non termina mai (capiterà, con i cicli infiniti del
        capitolo~\ref{chap:flusso}), è il modo per fermarlo.
  \item \tasti{Ctrl,D}, scritto come input a un programma, segnala la
        \textbf{fine dell'input}; lo useremo dal
        capitolo~\ref{chap:primi-programmi}.
  \item Nel terminale \tasti{Ctrl,C} non copia: per copiare e incollare,
        nel terminale di VS Code si usano il tasto destro del mouse o
        \tasti{Ctrl,Shift,C} e \tasti{Ctrl,Shift,V}.
\end{itemize}

\section{Il primo programma}

È il momento di mettere insieme i pezzi. Nella cartella
\code{programmazione} crea il file \code{ciao.c} e scrivici, esattamente
così, il programma seguente.

\begin{lstlisting}[style=c, name=ciao.c]
#include <stdio.h>

int main(void)
{
    printf("Ciao, mondo!\n");
    return 0;
}
\end{lstlisting}

Il significato di ogni riga sarà spiegato nel
capitolo~\ref{chap:primi-programmi}; per ora basta sapere che
\code{printf} stampa il testo fra virgolette, e che \code{\textbackslash n}
significa ``vai a capo''. Prima di proseguire, \textbf{salva} il file con
\tasti{Ctrl,S} (su Mac \tasti{Cmd,S}).

\begin{figure}[H]
  \centering
  \begin{tikzpicture}[font=\Lato\scriptsize]
    \foreach \x/\segno in {0/{\color{white}$\bullet$}, 7/{\color{dmdim}\texttimes}} {
      \fill[dmbar] (\x,0) rectangle ++(5,-0.5);
      \fill[dmbg] (\x,0) rectangle ++(2.3,-0.5);
      \fill[dmaccent] (\x,0) rectangle ++(2.3,-0.04);
      \node[anchor=west, text=white, font=\Lato\scriptsize] at (\x+0.2,-0.26)
        {{\FiraCodeNL\bfseries\color{vsciconblue}C}\ \ ciao.c\hspace{9pt}\segno};
    }
    \node[anchor=north, align=center] at (2.5,-0.65) {un pallino: modifiche \textbf{non salvate}};
    \node[anchor=north, align=center] at (9.5,-0.65) {una crocetta: il file è \textbf{salvato}};
  \end{tikzpicture}
  \caption{La scheda del file dice se ci sono modifiche non salvate.}
  \label{fig:salvato}
\end{figure}

gcc legge il file \textbf{dal disco}: se le modifiche non sono salvate,
compila la versione vecchia. È l'errore più frequente all'inizio
(``ho corretto ma non cambia niente!''). Per evitarlo una volta per tutte,
si può attivare il salvataggio automatico dal menu \menu{File,Auto Save}.

\subsection{Compilare ed eseguire}

Ora, nel terminale, compiliamo ed eseguiamo:

\begin{console}
$ gcc ciao.c
$ ls
a.out  ciao.c
$ ./a.out
Ciao, mondo!
\end{console}

\begin{itemize}
  \item \code{gcc ciao.c} compila il file. Non stampa nulla: nessun
        messaggio significa \textbf{nessun errore}. Il risultato è un nuovo
        file, l'eseguibile, che per tradizione si chiama \code{a.out}.
  \item \code{./a.out} esegue il programma, che stampa la sua riga.
\end{itemize}

Il nome \code{a.out} va bene per una prova, ma con più programmi nella
stessa cartella ogni compilazione sovrascriverebbe la precedente. Con
l'opzione \code{-o} (\textit{output}) si sceglie il nome dell'eseguibile; e
con \code{-Wall} (\textit{warnings: all}) si chiede a gcc di segnalare
anche le cose sospette, che non impediscono la compilazione ma spesso
nascondono un errore. È il modo in cui conviene compilare sempre:

\begin{console}
$ gcc -Wall -o ciao ciao.c
$ ./ciao
Ciao, mondo!
\end{console}

Ogni volta che si cambia il programma bisogna ripetere \textbf{tutto il
giro} (figura~\ref{fig:ciclo}): l'eseguibile è una traduzione fatta in un
certo momento, e non si aggiorna da solo quando cambia il testo.

\begin{figure}[H]
  \centering
  \begin{tikzpicture}[font=\Lato\footnotesize,
      fase/.style={wtnode, minimum width=3.0cm, minimum height=1.05cm}]
    \node[fase] (m) at (0,0) {\textbf{modifica}\\[1pt]{\scriptsize nell'editor}};
    \node[fase] (s) at (4.6,0) {\textbf{salva}\\[1pt]\tasti{Ctrl,S}};
    \node[fase, wtnodeacc] (c) at (9.2,0) {\textbf{compila}\\[1pt]{\FiraCodeNL\scriptsize gcc -Wall -o ciao ciao.c}};
    \node[fase, wtnodeok] (e) at (13.8,0) {\textbf{esegui}\\[1pt]{\FiraCodeNL\scriptsize ./ciao}};
    \draw[wtarrow] (m) -- (s);
    \draw[wtarrow] (s) -- (c);
    \draw[wtarrow] (c) -- (e);
    \draw[wtarrow] (e.south) -- ++(0,-0.55) -| node[wtlab, pos=0.25] {il risultato non è quello atteso? si ricomincia} (m.south);
    \draw[wtarrowacc] (c.north) -- ++(0,0.5) -| node[wtlab, pos=0.25] {errori di compilazione: si corregge il testo} (m.north);
  \end{tikzpicture}
  \caption{Il ciclo di lavoro: modifica, salva, compila, esegui.}
  \label{fig:ciclo}
\end{figure}

\subsection{Perché \code{./}?}

Perché per eseguire il programma si scrive \code{./ciao} e non solo
\code{ciao}, mentre per gcc basta \code{gcc}? Quando si scrive il nome di un
comando, la shell non lo cerca dappertutto: lo cerca solo in un elenco di
cartelle prestabilite, quelle dove stanno i programmi installati (come
\code{/usr/bin}, dove si trova \code{gcc}). L'elenco è una
\textit{variabile d'ambiente} chiamata \code{PATH}, e per sicurezza non
contiene la cartella corrente. Scrivendo \code{./ciao} si dà il
\textbf{percorso} del file invece del nome di un comando: ``il programma
\code{ciao} che sta \textit{qui}'' (\code{.} è la cartella corrente).

\begin{console}
$ ciao
ciao: command not found
$ ./ciao
Ciao, mondo!
\end{console}

\subsection{Gli errori di compilazione}

Che cosa succede se il programma contiene un errore? Togliamo il punto e
virgola alla fine della riga di \code{printf} e ricompiliamo:

\begin{affianco}[0.42][c]
\begin{lstlisting}[style=c, name=errore.c]
#include <stdio.h>

int main(void)
{
    printf("Ciao, mondo!\n")
    return 0;
}
\end{lstlisting}
\risultato
\begin{console}
$ gcc -Wall -o errore errore.c
errore.c: In function (*@‘@*)main(*@’@*):
errore.c:5:29: error: expected (*@‘@*);(*@’@*) before (*@‘@*)return(*@’@*)
    5 |     printf("Ciao, mondo!\n")
      |                             ^
      |                             ;
    6 |     return 0;
      |     ~~~~~~
\end{console}
\end{affianco}

Il messaggio va letto con calma, perché dice quasi tutto
(figura~\ref{fig:errore}): \textbf{dove} è l'errore (file, riga e colonna),
\textbf{di che tipo} è e \textbf{che cosa} si aspettava il compilatore. Le
righe successive ricopiano il codice e indicano con \code{\^{}} il punto
esatto; gcc arriva perfino a suggerire il \code{;} mancante. Quando c'è un
\code{error} l'eseguibile \textbf{non viene prodotto}: bisogna correggere e
ricompilare.

\begin{figure}[H]
  \centering
  \begin{tikzpicture}[font=\FiraCodeNL\normalsize,
      pezzo/.style={inner sep=0pt, anchor=base west, text=consolefg},
      eti/.style={font=\Lato\scriptsize, text=neutral!20!black, align=center}]
    \fill[consolebg, rounded corners=4pt] (-0.4,0.55) rectangle (11.6,-0.4);
    \node[pezzo] (f) at (0,0) {errore.c};
    \node[pezzo] (d1) at (f.base east) {:};
    \node[pezzo] (r) at (d1.base east) {5:29};
    \node[pezzo] (d3) at (r.base east) {:};
    \node[pezzo, text=red-gradient-2, font=\FiraCodeNL\normalsize\bfseries] (t) at ([xshift=6pt]d3.base east) {error:};
    \node[pezzo] (m) at ([xshift=6pt]t.base east) {expected ‘;’ before ‘return’};
    % sotto (punta in basso): file, riga e colonna, messaggio
    \foreach \p in {f, r, m}
      \draw[wtbrace] ([yshift=-0.55cm]\p.base west) -- ([yshift=-0.55cm]\p.base east);
    \node[eti, anchor=north] at ([yshift=-0.75cm]f.base) {il file};
    \node[eti, anchor=north] at ([yshift=-0.75cm]r.base) {riga 5,\\colonna 29};
    \node[eti, anchor=north] at ([yshift=-0.75cm]m.base) {che cosa non va: ``mi aspettavo\\un \FiraCodeNL;\Lato\ prima di \FiraCodeNL return\Lato''};
    % sopra (punta in alto): il tipo di messaggio
    \draw[wtbrace] ([yshift=0.75cm]t.base east) -- ([yshift=0.75cm]t.base west);
    \node[eti, anchor=south] at ([yshift=0.95cm]t.base) {il tipo: \textbf{error} (errore) o \textbf{warning} (avviso)};
  \end{tikzpicture}
  \caption{Anatomia di un messaggio di errore di gcc.}
  \label{fig:errore}
\end{figure}

Si noti che l'errore è segnalato alla riga \textbf{5}, alla fine della
riga di \code{printf}, anche se il compilatore se ne accorge solo leggendo
\code{return} alla riga 6: il punto e virgola poteva mancare solo lì. In
generale l'errore sta nella riga indicata \textbf{o poco prima}.

Esistono due tipi di messaggi:

\begin{itemize}
  \item un \code{error} è un errore vero: il testo non è C valido e
        l'eseguibile non viene creato;
  \item un \code{warning} (avviso) segnala qualcosa di lecito ma sospetto:
        l'eseguibile viene creato, ma quasi sempre il warning indica un
        errore di ragionamento. \textbf{I warning vanno trattati come
        errori}: un programma compilato con \code{-Wall} non deve
        produrne nessuno.
\end{itemize}

\info{Leggere solo il primo errore}{
  Un solo errore può confondere il compilatore e generarne a cascata molti
  altri, alcuni dei quali senza senso. La regola pratica: si corregge
  \textbf{il primo} messaggio dell'elenco, si ricompila, e si ripete.}

Con l'estensione C/C++ installata, VS Code si accorge di molti errori
\textbf{prima} della compilazione e li sottolinea con un'ondina rossa;
passando con il mouse sopra la sottolineatura compare la spiegazione
(figura~\ref{fig:ondina}). Il numero di errori compare anche nella barra di
stato, accanto a \icoerr[black].

\begin{figure}[H]
  \centering
  \begin{tikzpicture}[font=\FiraCodeNL\scriptsize,
      riga/.style={anchor=base west, inner sep=0pt},
      n/.style={anchor=base east, inner sep=0pt, text=vsclinenum, font=\FiraCodeNL\tiny}]
    \fill[dmbg, rounded corners=3pt] (0,0.35) rectangle (11.5,-1.95);
    \node[n] at (0.45,0) {4};
    \node[riga, text=vscfg] at (0.7,0) {\{};
    \node[n] at (0.45,-0.42) {5};
    \node[riga] at (0.7,-0.42) {\ \ \ \ {\color{vscfunc}printf}{\color{vscfg}(}{\color{vscstring}"Ciao, mondo!\textbackslash n"}{\color{vscfg})}};
    \node[n] at (0.45,-0.84) {6};
    \node[riga] (sp) at (0.7,-0.84) {\ \ \ \ };
    \node[riga, text=vscpurple] (ret) at (sp.base east) {return};
    \node[riga, text=vscfg] at (ret.base east) {\ 0;};
    \draw[red-gradient-3, ondina] ([yshift=-1.6pt]ret.base west) -- ([yshift=-1.6pt]ret.base east);
    \node[anchor=north west, fill=dmbar, draw=dmborder, text=dmfg, font=\Lato\scriptsize,
          inner sep=5pt, rounded corners=2pt, align=left] at ([yshift=-6pt]ret.base west)
      {expected a ';' \ \ {\color{dmdim}C/C++(65)}};
  \end{tikzpicture}
  \caption{L'estensione C/C++ sottolinea l'errore mentre si scrive.}
  \label{fig:ondina}
\end{figure}

\subsection{E il pulsante “Run”?}

L'estensione C/C++ aggiunge in alto a destra nell'editor un pulsante
\vscico[black]{gioca} (\textit{Run C/C++ File}) che compila ed esegue il
file aperto con un clic; la prima volta chiede quale compilatore usare (si
sceglie \textit{gcc}). È comodo, ma in questa dispensa compileremo
\textbf{sempre dal terminale}: è ciò che si fa a lezione e all'esame,
permette di scegliere le opzioni di gcc, e soprattutto fa vedere che cosa
succede davvero, passo per passo. Dal capitolo~\ref{chap:make} useremo
\code{make}, che si lancia anch'esso dal terminale.

\section{Quando qualcosa non funziona}

\begin{center}
\renewcommand{\arraystretch}{1.45}
\begin{tabular}{@{}>{\raggedright\arraybackslash}p{5.6cm}>{\raggedright\arraybackslash}p{10.6cm}@{}}
  \toprule
  \textbf{Sintomo} & \textbf{Causa probabile e rimedio} \\
  \midrule
  \code{gcc: command not found} &
    Il compilatore non è installato (ripeti il passo 2). Su Windows,
    controlla di essere nel terminale di \textbf{Ubuntu} (prompt con
    \code{\$}) e non in PowerShell (prompt \code{PS C:\textbackslash...}). \\
  \code{code: command not found} &
    Su Windows chiudi e riapri il terminale di Ubuntu dopo aver installato
    VS Code (il \code{PATH} si aggiorna solo all'apertura); su Mac esegui
    \textit{Shell Command: Install `code' command in PATH}. \\
  \codel{ciao.c: No such file or directory} &
    Il file non è nella cartella corrente: controlla con \code{pwd} e
    \code{ls} dove ti trovi e che cosa c'è, oppure il file non è stato
    salvato con quel nome. \\
  \codel{./ciao: No such file or directory} &
    Il programma non è stato compilato, oppure la compilazione è fallita
    (c'erano degli \code{error}), oppure l'eseguibile ha un altro nome
    (\code{a.out}?). \\
  Ho corretto ma il programma fa ancora la cosa sbagliata &
    Hai salvato? Hai ricompilato? Stai eseguendo l'eseguibile giusto? \\
  Il programma non termina &
    Probabilmente un ciclo infinito, oppure il programma aspetta un input
    da tastiera. \tasti{Ctrl,C} lo interrompe. \\
  In VS Code \code{\#include <stdio.h>} è sottolineato in rosso &
    L'estensione C/C++ non trova il compilatore. Su Windows: VS Code non
    è collegato a Ubuntu (guarda l'indicatore remoto) o l'estensione non è
    installata \textit{in WSL}. Altrove: verifica che \code{gcc -{}-version}
    funzioni. \\
  \bottomrule
\end{tabular}
\end{center}

\insintesi{
  Per scrivere ed eseguire un programma C servono tre strumenti: un
  \textbf{editor} (VS Code) per scrivere il testo del programma in un file
  \code{.c}, un \textbf{compilatore} (gcc) che lo traduce in linguaggio
  macchina, e un \textbf{terminale} da cui lanciare il compilatore e il
  programma. Su Windows tutto ciò vive in Ubuntu, grazie a WSL. Il ciclo di
  lavoro è sempre lo stesso: \textit{modifica, salva, compila, esegui}; e
  i messaggi di gcc indicano file, riga e colonna del problema.}

\section{Riferimenti}

\begin{itemize}
  \item \textbf{Installare Linux in Windows con WSL} --- Microsoft Learn. \\
        \urlx{https://learn.microsoft.com/it-it/windows/wsl/install}
  \item \textbf{Using C++ and WSL in VS Code} --- documentazione di VS Code
        (vale anche per il C). \\
        \urlx{https://code.visualstudio.com/docs/cpp/config-wsl}
  \item \textbf{Using Clang in Visual Studio Code} (macOS) e \textbf{Using
        C++ on Linux in VS Code} --- documentazione di VS Code. \\
        \urlx{https://code.visualstudio.com/docs/cpp/config-clang-mac} \\
        \urlx{https://code.visualstudio.com/docs/cpp/config-linux}
  \item \textbf{User Interface} --- la guida alle zone della finestra di
        VS Code. \\
        \urlx{https://code.visualstudio.com/docs/getstarted/userinterface}
\end{itemize}
